\documentclass[10pt,a4paper]{amsart}
\usepackage[margin=19mm]{geometry}
\usepackage{amsmath,amssymb,mathtools}
\usepackage[scr=rsfs,cal=euler]{mathalfa}
\usepackage{xfrac}
\usepackage[unicode,hidelinks,bookmarks=false]{hyperref}
\newcommand{\ie}{i.e.\ }
\newcommand{\cf}{cf.\ }
\newcommand{\resp}{respectively\ }
\newcommand{\Subsection}{\subsection}
\providecommand{\nobreakdash}{\nobreak\hbox{-}}
\setlength{\emergencystretch}{2em}
\multlinegap0pt
\usepackage{bm}
\usepackage{bbm}
\usepackage{dsfont}
\usepackage{xspace}
\usepackage{cancel}
\usepackage{mathtools}
\usepackage{amsthm}
\makeatletter
\@ifundefined{lemma}{\newtheorem{lemma}{Lemma}}{}
\makeatother
\newcommand\X{\mathrm{X}}
\newcommand\Z{\mathrm{Z}}
\newcommand\bbbeta{\boldsymbol{\beta}}
\newcommand\bomega{\boldsymbol{\omega}}
\newcommand\btheta{\boldsymbol{\theta}}
\newcommand\bu{\boldsymbol{u}}
\newcommand\bv{\boldsymbol{v}}
\newcommand\bx{\boldsymbol{x}}
\newcommand\curl{\mathop{\mathbf{curl}}\nolimits}
\newcommand\disp{\displaystyle}
\renewcommand\O{\Omega}
\title[Analysis and approximation of a vorticity-velocity-pressure ]{Analysis and approximation of a vorticity-velocity-pressure formulation for the Oseen equations}
\author{Veronica Anaya, Afaf Bouharguane, David Mora, Carlos Reales, Ricardo Ruiz Baier, Nour Seloula, Hector Torres}
\date{Source: arXiv:1805.01706v1 (CC BY 4.0)}
\begin{document}
\maketitle

\noindent\emph{Source and licence.} Analysis and approximation of a vorticity-velocity-pressure formulation for the Oseen equations. Version arXiv:1805.01706v1 (CC BY 4.0), distributed under the Creative Commons Attribution 4.0 International licence, as recorded in the arXiv licence metadata for this exact version and shown on the abstract page https://arxiv.org/abs/1805.01706v1. Full rights notice: licenses/arxiv-1805.01706-CC-BY-4.0.txt.\par\smallskip

\noindent\textit{Editorial scope.} Counted unit: the Lemma whose source label is \texttt{le1:continuous-problem} in the article (source0.tex lines 477--496, source label \texttt{le1:continuous-problem}), delivered together with the article's own complete original proof of it (source0.tex lines 498--543, one \texttt{proof} environment of 46 lines). No mathematics is written, repaired or re-derived: statement and proof are the article's own text line for line. Required context retained uncounted: none. Every definition, display and label of those lines is retained unchanged. Omitted from this package and not redistributed: 1--476, 497--497, 544--1878. Nothing inside a retained excerpt is omitted or altered: every retained body is delivered exactly as the article's lines are written, so no omission receipt of any kind accompanies this package. Citations: the retained excerpts carry no citation command at all, so no bibliography entry of the article is transcribed and no reference is redistributed. Reference adaptation: none was needed and none was made; every reference inside the delivered unit resolves inside it, so no unresolved reference or citation is produced. Printed slips kept literal: no printed word, symbol, hypothesis or number of the counted statement or of its proof was altered. Layout and class: the article's own class is \texttt{article} with packages including hyperref, amssymb, amsmath, epsfig, graphics, psfrag, graphicx, color, cite, fancyhdr, subfigure; the wrapper is the lane's pinned \texttt{amsart} editorial wrapper at 10pt on A4, which restates the theorem-like environments the retained text needs and adds the article's own macro definitions that the retained text needs (\texttt{\textbackslash O}, \texttt{\textbackslash X}, \texttt{\textbackslash Z}, \texttt{\textbackslash bbbeta}, \texttt{\textbackslash bomega}, \texttt{\textbackslash btheta}, \texttt{\textbackslash bu}, \texttt{\textbackslash bv}, \texttt{\textbackslash bx}, \texttt{\textbackslash curl}, \texttt{\textbackslash disp}), with extra wrapper packages (bm, bbm, dsfont, xspace, cancel, mathtools). The delivered numbering is the unit's own. Assets: no retained span contains a figure environment, a graphics-inclusion command, a PDF-page inclusion, a table or a TikZ picture; no image file is copied into this package, the package declares no asset dependency and no image file is redistributed with it. Licence: version arXiv:1805.01706v1 of this article is distributed under the Creative Commons Attribution 4.0 International licence, as recorded in the arXiv licence metadata for this exact version and shown on the abstract page https://arxiv.org/abs/1805.01706v1. A full rights notice is delivered at licenses/arxiv-1805.01706-CC-BY-4.0.txt. Characters: every retained excerpt is pure ASCII, so the delivered engine, XeLaTeX with its default Latin Modern text font, reports no missing character.
\begin{lemma}\label{le1:continuous-problem}
Let us assume that 
\begin{equation}\label{beta-assumption}
\frac{2\Vert\bbbeta\Vert_{\infty,\O}^2}{\nu\sigma}< 1,
\end{equation}
and let us define the bilinear form $\mathcal{A}(\cdot,\cdot)$ specified as 
$$\mathcal A((\bu,\bomega),(\bv,\btheta)):=a(\bu,\bv)+\;b_{1}(\bv,\bomega)+b_{1}(\bu,\btheta)-\;d(\bomega,\btheta)+c(\bomega,\bv).$$
Then, there exist $\alpha_1,\alpha_2>0$ such that
\begin{equation}\label{acot}
\vert\mathcal A((\bu,\bomega),(\bv,\btheta))\vert\le \alpha_1\Vert(\bu,\bomega)\Vert_{\mathcal X}\Vert(\bv,\btheta)\Vert_{\mathcal X},
\end{equation}
and
\begin{equation}\label{cont-inf-sup-Aa}
\disp\,\sup_{\stackrel{\scriptstyle (\bv,\btheta)\in {\mathcal X}}{(\bv,\btheta)\ne0}}
\frac{\mathcal A((\bu,\bomega),(\bv,\btheta))}{\|(\bv,\btheta)\|_{\mathcal X}}\,\ge\,\alpha_2\,\|(\bu,\bomega)\|_{\mathcal X}
\quad\forall\,(\bu,\bomega)\in \mathcal X,
\end{equation}
where $\mathcal X:=\X\times\Z$, endowed
with the corresponding product norm, is a Hilbert space.
\end{lemma}

\begin{proof}
As a consequence of the boundedness of $a(\cdot,\cdot)$
$b_1(\cdot,\cdot)$, $c(\cdot,\cdot)$, and $d(\cdot,\cdot)$,  the 
bilinear form $\mathcal A(\cdot,\cdot)$ is bounded and so condition \eqref{acot} readily 
follows.

Concerning the satisfaction of the inf-sup condition \eqref{cont-inf-sup-Aa}, for a 
 given $(\bu,\bomega)\in\mathcal X$, we can define 
$$\tilde{\btheta}:=-\bomega\in\Z,\quad \text{and}\quad \tilde{\bv}:=(\bu+\hat{c}\sqrt{\nu}\curl\bomega)\in\X,$$
where $\hat{c}>0$ is a constant to be chosen later. 
%
We can then immediately assert that 
\begin{align*}
\mathcal A((\bu,\bomega),(\tilde \bv,\tilde\btheta))
& =\sigma\int_{\O}\bu\cdot\tilde{\bv}\, d\bx+\sqrt{\nu}\int_{\O}\curl\bomega\cdot\tilde{\bv}\, d\bx
+\sqrt{\nu}\int_{\O}\curl\tilde{\btheta}\cdot\bu\, d\bx \\
&\quad -\int_{\O}\bomega\cdot\tilde{\btheta}\, d\bx+\frac{1}{\sqrt{\nu}}\int_{\O}(\bomega\times\bbbeta)\cdot\tilde{\bv}\, d\bx\\
&\geq \sigma\Vert\bu\Vert_{0,\O}^2+\hat{c}\sqrt{\nu}\sigma\int_{\O}\bu\cdot\curl\bomega\, d\bx
+\sqrt{\nu}\int_{\O}\bu\cdot\curl\bomega\, d\bx+\hat{c}\nu\Vert\curl\bomega\Vert_{0,\O}^2 \\
&\quad -\sqrt{\nu}\int_{\O}\bu\cdot\curl\bomega\, d\bx+\Vert\bomega\Vert_{0,\O}^2+\frac{1}{\sqrt{\nu}}\int_{\O}(\bomega\times\bbbeta)\cdot\bu\, d\bx
+\hat{c}\int_{\O}(\bomega\times\bbbeta)\cdot\curl\bomega\, d\bx\\
& \geq \sigma\Vert\bu\Vert_{0,\O}^2-\frac{\sigma}{4}\Vert\bu\Vert_{0,\O}^2
-\hat{c}^2\sigma\nu\Vert\curl\bomega\Vert_{0,\O}^2
+\hat{c}\nu\Vert\curl\bomega\Vert_{0,\O}^2+\Vert\bomega\Vert_{0,\O}^2
-\frac{2\sigma}{3}\Vert\bu\Vert_{0,\O}^2\\
&\quad -\frac{3\Vert\bbbeta\Vert_{\infty,\O}^2}{2\nu\sigma}\Vert\bomega\Vert_{0,\O}^2
-\frac{\Vert\bbbeta\Vert_{\infty,\O}^2}{2\nu\sigma}\Vert\bomega\Vert_{0,\O}^2
-2\hat{c}^2\sigma\nu\Vert\curl\bomega\Vert_{0,\O}^2\\
& =\frac{\sigma}{12}\Vert\bu\Vert_{0,\O}^2+\hat{c}\left(1-3\hat{c}\sigma\right)\nu\Vert\curl\bomega\Vert_{0,\O}^2
+\left(1-\frac{2\Vert\bbbeta\Vert_{\infty,\O}^2}{\nu\sigma}\right)\Vert\bomega\Vert_{0,\O}^2,
\end{align*}
where we have used the bound $\Vert\bomega\times\bbbeta\Vert_{0,\O}\le2\Vert\bbbeta\Vert_{\infty,\O}\Vert\bomega\Vert_{0,\O}$.
Choosing $\hat{c}=1/(4\sigma)$ and exploiting \eqref{beta-assumption}, we arrive at 
$$\mathcal A((\bu,\bomega),(\tilde \bv,\tilde\btheta))\ge C\Vert(\bu,\bomega)\Vert_{\mathcal X}^2,$$
with $C$ independent of $\nu$. 
%
On the other hand, by construction we realise that   
$\Vert{\tilde{\btheta}}\Vert_{\Z}=\Vert\bomega\Vert_{\Z}$ and 
$\Vert\tilde{\bv}\Vert_{0,\O}\le C\hat{c}(\Vert\bu\Vert_{0,\O}+\Vert\bomega\Vert_{\Z})$,
and consequently 
$$\sup_{\stackrel{\scriptstyle (\bv,\btheta)\in\mathcal X}{(\bv,\btheta)\ne0}}\frac{\mathcal A((\bu,\bomega),(\bv,\btheta))}
{\Vert(\bv,\btheta)\Vert_{\mathcal X}}\ge \frac{\mathcal A((\bu,\bomega),(\tilde \bv,\tilde\btheta))}
{\Vert(\tilde\bv,\tilde\btheta)\Vert_{\mathcal X}}
\ge \alpha_2\Vert(\bu,\bomega)\Vert_{\mathcal X}\qquad\forall(\bu,\bomega)\in\mathcal X,$$
which finishes the proof.
\end{proof}

\end{document}
